\documentclass[12pt,aspectratio=169]{beamer} % 16:9 para slides largos

% -------------------------------------------------
% Tema e pacotes
% -------------------------------------------------
\usetheme{Madrid}
\usecolortheme{dolphin} % azul elegante
\usepackage[utf8]{inputenc}
\usepackage[brazil]{babel}
\usepackage{graphicx}
\usepackage{hyperref}
\usepackage{float}

% -------------------------------------------------
% Reduzir margens para ocupar mais espaço
% -------------------------------------------------
\setbeamersize{text margin left=5mm, text margin right=5mm}

% -------------------------------------------------
% Título
% -------------------------------------------------
\title{Tema 3 – Introdução aos Softwares de Revisão Sistemática}
\subtitle{Apontamento de Literatura – Revisão Sistemática e Meta-análises}
\author{Pós-graduação (Mestrado e Doutoramento)}
\date{}

\begin{document}

\frame{\titlepage}

% =================================================
\begin{frame}{Conceito de software de revisão sistemática}
Softwares de revisão sistemática são ferramentas digitais projetadas para \textbf{organizar, gerenciar, avaliar e sintetizar a literatura científica} de forma rigorosa e reprodutível.

\textbf{Objetivos principais:}
\begin{itemize}
    \item Facilitar busca, seleção e triagem de estudos;
    \item Automatizar extração de dados e avaliação de qualidade;
    \item Auxiliar na síntese qualitativa e quantitativa;
    \item Garantir transparência e rastreabilidade.
\end{itemize}

\begin{figure}[H]
    \centering
    \includegraphics[width=0.9\textwidth]{fig1_softwares.png} % maior largura
    \caption{Diagrama conceptual de softwares de revisão sistemática}
\end{figure}
\end{frame}

% =================================================
\begin{frame}{Softwares mais utilizados}
\begin{table}[]
\centering
\begin{tabular}{|p{3cm}|p{3cm}|p{9cm}|} % largura maximizada
\hline
\textbf{Software} & \textbf{Tipo} & \textbf{Finalidade / Destaques} \\ \hline
EPPI-Reviewer Web & Plataforma online & Triagem, codificação, meta-análise, projetos personalizáveis \\ \hline
JBI SUMARI & Desktop & Revisões sistemáticas e de escopo, avaliação crítica de estudos \\ \hline
RevMan & Desktop (Cochrane) & Gestão de revisões Cochrane, metanálises, gráficos de efeito \\ \hline
Rayyan & Online & Triagem rápida, detecção de duplicados, colaboração entre revisores \\ \hline
Outros & Variável & Mendeley, Zotero, Covidence, DistillerSR, etc. \\ \hline
\end{tabular}
\caption{Principais softwares de revisão sistemática}
\end{table}
\end{frame}

% =================================================
\begin{frame}{EPPI-Reviewer Web}
EPPI-Reviewer é uma plataforma online robusta para gerenciar grandes revisões sistemáticas.

\textbf{Principais funcionalidades:}
\begin{itemize}
    \item Projetos colaborativos;
    \item Importação de referências de bases de dados (PubMed, Scopus, Web of Science);
    \item Triagem e codificação estruturada;
    \item Ferramentas avançadas de síntese e metanálise;
    \item Exportação de relatórios padronizados.
\end{itemize}

\begin{figure}[H]
    \centering
    \includegraphics[width=0.85\textwidth]{fig2_eppi.png} % maior largura
    \caption{Interface do EPPI-Reviewer Web}
\end{figure}
\end{frame}

% =================================================
\begin{frame}{Review Manager (RevMan)}
RevMan é o software oficial da Cochrane para revisões sistemáticas e metanálises.

\textbf{Características principais:}
\begin{itemize}
    \item Desenvolvimento de protocolos e revisões completas;
    \item Organização de referências e estudos incluídos;
    \item Meta-análise com gráficos (forest plots, funnel plots);
    \item Avaliação de risco de viés;
    \item Exportação para publicação.
\end{itemize}

\begin{figure}[H]
    \centering
    \includegraphics[width=0.85\textwidth]{fig3_revman.png} % maior largura
    \caption{Interface do RevMan para gestão de revisões sistemáticas}
\end{figure}
\end{frame}

% =================================================
\begin{frame}{Funcionalidades detalhadas do RevMan}
\begin{table}[]
\centering
\begin{tabular}{|p{6cm}|p{10cm}|} % largura maximizada
\hline
\textbf{Função} & \textbf{Descrição} \\ \hline
Inserção de estudos & Adicionar e organizar estudos incluídos e excluídos \\ \hline
Avaliação de qualidade & Risk of Bias tool; checklist metodológica \\ \hline
Meta-análise & Cálculo automático de efeitos, forest plots \\ \hline
Gráficos de síntese & Funnel plots, gráficos de heterogeneidade \\ \hline
Exportação & PDF, Word e dados para bases externas \\ \hline
\end{tabular}
\caption{Principais funcionalidades do RevMan}
\end{table}
\end{frame}

% =================================================
\begin{frame}{Cochrane e RevMan}
\begin{itemize}
    \item Cochrane: organização internacional para revisões sistemáticas de alta qualidade;
    \item RevMan é gratuito para membros Cochrane; versão básica gratuita para não membros;
    \item Objetivo: padronizar revisões, aumentar confiabilidade e facilitar disseminação da evidência.
\end{itemize}

\begin{figure}[H]
    \centering
    \includegraphics[width=0.85\textwidth]{fig4_cochrane.png} % maior largura
    \caption{Fluxo de produção de revisões sistemáticas na Cochrane}
\end{figure}
\end{frame}

% =================================================
\begin{frame}{O que o RevMan pode fazer por si}
\begin{itemize}
    \item Automatiza cálculos estatísticos de metanálise;
    \item Produz gráficos claros de efeito e heterogeneidade;
    \item Permite organização de referências e estudos incluídos;
    \item Ajuda a preparar manuscritos padronizados;
    \item Facilita colaboração em equipe e rastreabilidade.
\end{itemize}

\begin{figure}[H]
    \centering
    \includegraphics[width=0.85\textwidth]{fig5_revman_func.png} % maior largura
    \caption{Exemplo de funcionalidades do RevMan}
\end{figure}
\end{frame}

% =================================================
\begin{frame}{Considerações finais}
Para mestrados e doutoramentos:
\begin{itemize}
    \item Domínio de softwares de revisão sistemática é essencial;
    \item EPPI-Reviewer, RevMan e Rayyan oferecem abordagens complementares;
    \item Escolha depende de tipo de revisão, tamanho da amostra e recursos disponíveis;
    \item Saber explorar metanálises, gráficos e relatórios automáticos aumenta rigor e eficiência.
\end{itemize}
\end{frame}

\end{document}
